\section{Rangkuman}
\begin{itemize}
    \item \textit{Block cipher} membagi plainteks menjadi blok tetap $n$ bit ($C=E_K(P)$, $P=D_K(C)$); panjang cipherteks sama dengan panjang plainteks, sedangkan panjang kunci eksternal bebas.
    \item Untuk setiap kunci, $E_K$ harus merupakan permutasi bijektif atas himpunan $2^n$ blok; cipher blok hanya memilih $2^m$ permutasi dari $2^n!$ permutasi yang tersedia.
    \item Bila pesan tidak habis dibagi ukuran blok, blok terakhir ditambahi bit \textit{padding} agar panjangnya kembali $n$ bit.
    \item Confusion menyembunyikan hubungan kunci--cipherteks melalui substitusi non-linear (S-box, misalnya $S_1$ DES: $\code{110100}\to\code{1001}$), sedangkan diffusion menyebarkan pengaruh satu bit plainteks ke banyak bit cipherteks melalui permutasi.
    \item Confusion dan diffusion baru tercapai menyeluruh setelah substitusi dan permutasi diulang berkali-kali sebagai \textit{iterated cipher}, dengan kunci putaran yang berbeda pada setiap putaran.
    \item Mode ECB ($C_i=E_K(P_i)$) sederhana dan mendukung akses acak, tetapi blok plainteks yang sama selalu menghasilkan cipherteks yang sama sehingga pola data bocor dan blok dapat dibuang atau disusun ulang.
    \item Mode CBC ($C_i=E_K(P_i\oplus C_{i-1})$, $C_0=IV$) menghapus pola berulang dan menutup manipulasi blok, namun galat satu bit menjalar ke seluruh blok cipherteks berikutnya saat enkripsi.
    \item Mode CFB ($C_j=P_j\oplus\mathrm{MSB}_s(E_K(X_j))$) memperlakukan cipher blok seperti cipher alir untuk data yang belum genap satu blok, dengan umpan balik diambil dari cipherteks.
    \item Mode OFB mengambil umpan balik dari keluaran $E_K$, bukan dari cipherteks, sehingga aliran kuncinya tidak bergantung pada data dan galat tidak menjalar sama sekali.
    \item Mode CTR ($C_i=P_i\oplus E_K(\text{counter}_i)$) mendukung pemrosesan paralel penuh dan akses acak karena tiap blok tidak bergantung pada blok lain; nilai counter wajib tidak pernah berulang.
    \item Jaringan Feistel ($L_i=R_{i-1}$, $R_i=L_{i-1}\oplus f(R_{i-1},K_i)$) bersifat reversibel sehingga dekripsi memakai urutan kunci terbalik tanpa fungsi $f^{-1}$ terpisah; sifat ini tidak bergantung pada bentuk fungsi $f$.
    \item Jaringan SPN seperti AES menuntut setiap komponen putaran dapat dibalik, sehingga inversi S-box dan matriks pencampur harus dirancang khusus.
    \item \textit{Key scheduling} menurunkan kunci eksternal menjadi $r$ kunci putaran yang berbeda; DES mengubah kunci 64 bit menjadi 16 kunci putaran 48 bit.
    \item Ukuran blok menentukan luas diffusion, sedangkan ukuran kunci menentukan ketahanan terhadap pencarian menyeluruh; blok 64 bit perlu diganti kunci sebelum $2^{32}$ blok karena paradoks ulang tahun.
\end{itemize}
